\documentclass[11pt]{article}
\usepackage[a4paper,margin=25mm]{geometry}
\usepackage{amsmath,amssymb}
\usepackage[T1]{fontenc}
\usepackage{lmodern}
\usepackage{hyperref}
\hypersetup{colorlinks=true,urlcolor=blue}
\setlength{\parindent}{0pt}
\setlength{\parskip}{7pt}
\title{Four colors suffice for all prime distances on the real line}
\author{SucRunBug}
\date{1 October 2026}
\begin{document}
\maketitle
\textbf{Conjecture 00000000024 -- Verdict: FALSE.}

\textbf{Filed statement.} The prime-distance graph on $\mathbb R$ has chromatic number $\aleph_0$, and the chromatic numbers for finite initial prime-distance sets tend to infinity.

\section*{Proof}
Define adjacency by $|x-y|=p$ for some positive integer prime $p$. Consider the coloring
\[c(x)=\lfloor x\rfloor\pmod 4,\qquad x\in\mathbb R.\]
Its four residues are the colors. If $y=x+p$, the integer-translation identity for the floor function gives
\[\lfloor y\rfloor=\lfloor x\rfloor+p.\]
No prime is divisible by $4$: the only even prime is $2$, and $4$ does not divide $2$. Consequently $c(x)\ne c(y)$. The case $x=y+p$ follows by symmetry. Thus this is a proper coloring of the graph on \emph{all} real numbers, including negative and nonintegral numbers.

The same coloring is proper after restricting the allowed prime distances to any subset. In particular, all finite initial prime-distance graphs have chromatic number at most $4$. They therefore cannot have chromatic numbers tending to infinity, and the full graph cannot have chromatic number $\aleph_0$.

For completeness, the bound is sharp. The four real points $0,2,5,7$ have pairwise distances $2,5,7,3,5,2$, all prime. They form a $K_4$, so the full chromatic number is exactly $4$. The finite initial prime-distance graphs also have chromatic number exactly $4$ once the distance $7$ is included.

\section*{Lean formalization}
The Lean project defines the actual simple graph on $\mathbb R$, constructs a proper coloring with color type \texttt{Fin 4}, and proves the same result for every subset of prime distances. The supplementary $K_4$ argument establishes sharpness in the written proof; the refutation needs only the fully formalized four-color upper bound.

\textbf{Reproduction.} In the supplied \texttt{lean4/} project, run
\texttt{lake exe cache get}, then \texttt{lake build} and
\texttt{lake env lean Check.lean}. The pinned versions are Lean and Mathlib
\texttt{v4.33.0}. The supplied \texttt{reproduce.py} performs independent
finite sanity checks; these checks are not the proof of any infinite assertion.

\textbf{Source.} \url{https://github.com/The-Last-Math-Competition/The-Last-Math-Competition/blob/main/conjectures/00000000024.md}
\end{document}
